\section{Предлагаемый подход и реализация прототипа}
\label{sec:Chapter4} \index{Chapter4}

Как было показано в предыдущем разделе, основная неэффективность существующих сканеров уязвимостей при работе с container registry связана с полной загрузкой слоев, в то время как для анализа в большинстве случаев требуется лишь небольшая часть их содержимого. В данном разделе обсуждаются возможные подходы к решению этой проблемы, обосновывается выбор формата eStargz и описывается реализация прототипа.

\subsection{Возможные направления решения}

Возможные направления решения сформулированной проблемы можно разделить на две большие категории:

\begin{itemize}
\item \textit{внеcение изменений на сторону registry}~--- например, добавление специального API, позволяющего запрашивать произвольные файлы из содержимого слоев. Этот путь неэффективен, так как для обеспечения гранулярной загрузки файлов все равно потребуется распаковка образов на стороне сервера, что создает дополнительную нагрузку на инфраструктуру и требует существенных изменений относительно спецификации OCI;
\item \textit{использование специальных форматов слоев, поддерживающих частичную загрузку}~--- содержимое слоев огранизуется так, чтобы клиент мог запрашивать произвольные файлы, не загружая слой целиком. К этой категории относятся форматы eStargz, RAFS (используемый в Nydus), и др. На момент написания работы этот путь активно применяется в области ускорения запуска контейнеров.
\end{itemize}

Так, в рамках данной работы был сделан упор на изучение форматов слоев с поддержкой ленивой загрузки.

\subsection{Обзор форматов ленивой загрузки}

\subsubsection{Предпосылки}

Идея ленивой загрузки слоев в контексте контейнерных образов восходит к работе Slacker~\cite{harter2016slacker}. Авторы Slacker исследовали профиль работы Docker и обнаружили, что загрузка образов занимает 76\% времени запуска контейнера, при этом только 6.4\% от загруженных данных в конечном счете используются. Это наблюдение легло в основу идеи запуска контейнера до полной загрузки его образа: содержимое подгружается лениво, по мере необходимости. Реализация Slacker использовала NFS в качестве хранилища, что ограничивало её применимость, однако сама идея ленивой загрузки получила в дальнейшем широкое развитие.

\subsubsection{eStargz}

\textit{eStargz}~\cite{estargzspec}~--- формат ленивой загрузки слоев, развиваемый в рамках проекта containerd. Он является развитием формата \textit{stargz}, изначально предложенного в проекте Google CRFS~\cite{crfs}, и впоследствии дополненного механизмами проверки целостности и оптимизации производительности.

Ключевая идея формата eStargz заключается в том, что вместо одного непрерывного gzip-потока, содержащего весь tar-архив, формат разбивает архив на отдельные gzip-потоки, каждый из которых содержит независимый <<кусочек>> содержимого. Возможность гранулярной загрузки файлов обеспечивается добавлением к архиву \textit{TOC (Table of Contents)}, содержащей описание всех файлов с указанием их позиции в архиве. Фактически TOC представляет из себя отдельный JSON-документ \texttt{stargz.index.json}, содержащийся в архиве. Кроме того, в конец архива добавляется специальная запись фиксированного размера, содержащая адрес начала TOC в архиве (т.н. \textit{footer}).

Структура архива eStargz проиллюстрирована рисунке \ref{fig:estargz}.

\begin{figure}[H]
    \centering
    \includegraphics[width=0.8\textwidth]{estargz.jpg}
    \caption{Структура архива eStargz}
    \label{fig:estargz}
\end{figure}

Принципиально важно, что формат eStargz \textit{полностью совместим} с обычным форматом слоев OCI: с точки зрения registry и стандартных средств распаковки eStargz-слой представляет собой обычный tar.gz-архив. Это означает, что:

\begin{itemize}
\item eStargz-слои можно загружать и хранить в любых OCI-совместимых registry без необходимости каких-либо изменений на стороне сервера;
\item eStargz-образы можно запускать обычными container runtime, не поддерживающими ленивую загрузку;
\end{itemize}

Поддержка ленивой загрузки на стороне клиента, в свою очередь, выглядит следующим образом:

\begin{itemize}
\item клиент скачивает footer слоя (фиксированного небольшого размера) с помощью HTTP Range-запроса~\cite{httprange};
\item из footer извлекается адрес TOC, по которому ещё одним Range-запросом скачивается сам TOC;
\item для каждого нужного клиенту файла из TOC извлекается адрес соответствующего gzip-потока, который запрашивается отдельным Range-запросом;
\item полученный gzip-поток разжимается и возвращает содержимое файла.
\end{itemize}

Размер footer eStargz составляет всего 51 байт, размер TOC~--- как правило, от десятков до сотен килобайт в зависимости от количества файлов в слое. Таким образом, для получения информации о содержимом слоя достаточно нескольких HTTP-запросов суммарным объемом в сотни килобайт; затем каждый нужный файл скачивается отдельным запросом размером в несколько килобайт.

\subsubsection{Альтернативы}

Помимо eStargz, в индустрии получили распространение и другие форматы ленивой загрузки:

\begin{itemize}
\item \textit{Nydus / RAFS}~\cite{nydus}~--- формат, разработанный в рамках проекта Dragonfly от Alibaba. Использует собственный формат RAFS, в котором метаданные и данные слоя хранятся раздельно. Поддерживает дедупликацию данных, в том числе между различными слоями различных образов, что позволяет дополнительно сэкономить трафик при загрузке. Однако RAFS слои \textit{не совместимы} с OCI и не могут быть использованы в container runtime без поддержки формата.
\item \textit{Seekable OCI (SOCI)}~\cite{soci}~--- технология, разработанная AWS. В отличие от eStargz и Nydus, SOCI не требует изменения самих слоев образа для обеспечения поддержки ленивой загрузки, что естественным образом обеспечивает полную обратную совместимость. Вместо этого создается отдельный артефакт~--- \textit{SOCI-индекс}. Такой индекс содержит TOC архива, а также \textit{набор контрольных точек состояния алгоритма сжатия} (т.н. \textit{zInfo}), которые позволяют разжимать не весь gzip-поток целиком, а лишь необходимый фрагмент между двумя контрольными точками (\textit{т.н. span}). В свою очередь, в TOC хранится информация не только о смещении в распакованном tar-архиве, но и о позиции ближайшего span'а и соответствующей ему контрольной точке.
\end{itemize}

Сравнительная таблица перечисленных форматов представлена в табл.~\ref{tab:lazy-pull-formats}.

\begin{table}[H]
    \centering
    \caption{Сравнение форматов ленивой загрузки слоев}
    \label{tab:lazy-pull-formats}
    \begin{tabular}{lccc}
        \toprule
        Формат & Совместимость с OCI & Дедупликация & Конвертация \\
        \midrule
        eStargz   & полная             & нет (на уровне файлов) & да (отдельный шаг) \\
        Nydus / RAFS & нет             & да (на уровне chunk)   & да (отдельный шаг) \\
        Seekable OCI & полная (отдельный индекс) & нет & да (генерация индекса) \\
        \bottomrule
    \end{tabular}
\end{table}

Из приведенного сравнения видно, что eStargz является наиболее <<щадящим>> с точки зрения интеграции в существующую инфраструктуру: он не требует изменений на стороне registry и не нарушает совместимость с обычными средствами обработки образов. С точки зрения адаптации в сканере уязвимостей это является решающим преимуществом, поскольку позволяет внедрить ленивую загрузку без выноса этой функциональности на сторону инфраструктуры. По этой причине именно eStargz выбран в качестве основы для дальнейшей работы.

\subsection{Ленивая загрузка слоёв через eStargz как метод минимизации трафика}

С учётом вышеизложенного, предлагаемый подход к решению поставленной проблемы заключается в следующем:

\begin{enumerate}
\item образы загружаются в registry в eStargz-формате (например, путем конвертации в рамках CI-процесса сразу после сборки, либо настройкой сборщика на немедленное создание eStargz);
\item сканер уязвимостей при анализе образа использует доступ к содержимому слоев через TOC eStargz, скачивая только те файлы, которые требуются включенным анализаторам;
\item остальная функциональность сканера~--- анализаторы, базы уязвимостей, форматы отчетов~--- остается без изменений.
\end{enumerate}

Такой подход обладает следующими преимуществами:

\begin{itemize}
\item не требует изменений в спецификации OCI и на стороне registry;
\item использует существующую инфраструктуру конвертации и поддержки eStargz, развивающуюся в рамках проекта containerd;
\item не привязан к конкретному сканеру~--- предложенный подход применим к любому сканеру, реализующему модель анализа слоев на основе фильтрации файлов.
\end{itemize}

Чтобы продемонстрировать применимость этого подхода на практике, в рамках работы был разработан прототип, реализующий ленивую загрузку слоев через eStargz в сканере Trivy. Описанию реализации прототипа посвящена оставшаяся часть данного раздела.

\subsection{Реализация ленивой загрузки в Trivy}

В качестве основы для разработки прототипа был выбран сканер Trivy~\cite{trivy} по совокупности следующих соображений:

\begin{itemize}
\item Trivy является наиболее популярным open-source сканером уязвимостей, и его адаптация имеет наибольшую практическую ценность;
\item модульная архитектура Trivy на основе анализаторов хорошо сочетается с моделью ленивой загрузки: можно использовать тот же механизм определения нужных файлов, что и в обычном режиме работы.
\end{itemize}

\subsubsection{Архитектура Trivy}

Прежде чем переходить к описанию вносимых изменений, имеет смысл коротко рассмотреть устройство Trivy в той части, которая для них важна. Внутренне Trivy организован модульно: центральной концепцией являются \textit{анализаторы (analyzers)}~--- независимые модули, каждый из которых реагирует на определенные типы файлов и извлекает из них необходимую информацию. Все анализаторы перечислены в пакете \texttt{pkg/fanal/analyzer}, при этом каждый анализатор имеет:

\begin{itemize}
\item метод \texttt{Required}, возвращающий, является ли данный файл интересным для анализатора;
\item метод \texttt{Analyze}, осуществляющий собственно анализ файла и возвращающий извлеченную из него информацию.
\end{itemize}

Помимо анализаторов, Trivy поддерживает \textit{пост-анализаторы (post-analyzers)}, работающие на наборе уже подобранных файлов целиком (например, для разрешения зависимостей по содержимому нескольких связанных манифестов). Совокупная логика сводится к тому, что в процессе обхода файловой системы образа Trivy для каждого файла опрашивает все включенные анализаторы; если хотя бы один из них объявил файл нужным, файл передается на анализ.

Дополнительно в коде Trivy имеется интерфейс \texttt{StaticPathAnalyzer}, который реализуется анализаторами с заранее известным фиксированным набором путей (apk, dpkg, rpm). Соответствующий метод \texttt{PartialStaticPaths()} возвращает совокупный список таких путей и используется в режимах, когда сканер может частично загружать содержимое образа (например, при работе с локальной файловой системой). Однако языковые анализаторы (pip, npm, go, jar и др.) этот интерфейс не реализуют, поскольку их интересуют файлы с переменным расположением, выявляемые методом \texttt{Required}. По этой причине \texttt{PartialStaticPaths} оказывается неполноценным критерием выбора файлов, а в режиме работы с registry этот механизм и вовсе не используется, и слои скачиваются целиком.

Для поддержки ленивой загрузки в архитектуру сканера потребовалось внести ряд изменений, не затрагивающих при этом основную логику работы анализаторов и сопоставления с базой уязвимостей.

\subsubsection{Пофайловая загрузка из registry}

В первую очередь был реализован механизм запроса произвольных байтовых диапазонов из registry поверх стандартных HTTP Range-запросов. Этот механизм позволяет сканеру обращаться к удаленному слою как к локальному файлу, запрашивая только нужные фрагменты данных. При этом учитывается необходимость аутентификации: если registry требует авторизации, сканер автоматически получает и использует соответствующие токены.

\subsubsection{Обход eStargz-слоя через TOC}

Центральным элементом ленивой загрузки является новый алгоритм обхода содержимого слоя. Вместо того чтобы скачивать слой целиком и распаковывать его, сканер выполняет следующие шаги:

\begin{enumerate}
\item скачивает небольшой фрагмент в конце слоя (footer), чтобы узнать расположение TOC;
\item скачивает сам TOC, представляющий собой легковесный JSON-документ со списком всех файлов в слое и их смещениями;
\item для каждого файла из TOC проверяет, нужен ли он для анализа;
\item если файл нужен, сканер вычисляет его точное расположение в сжатом слое и скачивает только этот фрагмент отдельным запросом.
\end{enumerate}

Важной оптимизацией на этом этапе стало чтение каждого нужного файла целиком за один запрос. Как показала практика, чтение файла мелкими порциями приводит к многократному дублированию HTTP-запросов к одному и тому же фрагменту архива, что сводит на нет всю экономию трафика.

\subsubsection{Фильтрация файлов для анализа}

Чтобы воспользоваться существующей системой анализаторов Trivy, необходимо было предоставить им возможность отвечать на вопрос <<нужен ли вам этот файл?>> без фактического скачивания его содержимого. Как было отмечено выше, существующий в Trivy механизм \texttt{PartialStaticPaths} для этой задачи не подходит, поскольку учитывает только системные анализаторы с фиксированным набором путей и пропускает языковые анализаторы~--- если опираться только на него, сканер не обнаруживал бы уязвимости в языковых пакетах.

Для решения этой проблемы был реализован новый, универсальный механизм фильтрации (метод \texttt{IsRequired}). Он последовательно опрашивает все включенные в сканере анализаторы и пост-анализаторы, передавая им имя файла и его метаданные из TOC слоя. Если хотя бы один анализатор (будь то системный или языковой) подтверждает, что файл ему интересен, сканер принимает решение о его загрузке. Таким образом, все анализаторы участвуют в выборке файлов, и ни одна потенциальная уязвимость не пропускается.

\subsubsection{Интеграция в конвейер сканирования}

Новый алгоритм обхода был интегрирован в существующий конвейер сканирования Trivy в качестве альтернативного пути обработки слоя. При включении соответствующей опции сканер обрабатывает слой в режиме ленивой загрузки. Принципиально важно, что после скачивания нужных файлов их содержимое передается в те же самые функции анализа, что и при обычном сканировании. Это гарантирует, что результаты сканирования в режиме ленивой загрузки будут абсолютно идентичны результатам обычного сканирования: меняется лишь способ доставки данных, а сама логика поиска уязвимостей остается неизменной.

\newpage
